% !TeX root = ../../surface-lagrangian-torus-fibration.tex
% 2026-11-21 / v1.3: 全 stage に共通な reduced marking と source-inclusive な飽和被覆．
\section{Source を含む neck と三重点 end の共通 exhaustion}
\label{ke:neck-corner-exhaustion-section}

固定残余基底の定理\ref{ke:projected-source-neck}と，
moving arm の命題\ref{ke:arm-common-moment-strip}を接続する．
選択の順序は，残余 compact 領域，normal cutoff，moment interval，
最後に退化 parameter の閾値である．この順序によって，固定 compact の定理を
未証明の moving-base 一様評価として用いることを避ける．
最初に source と saddle を含む neck の end germs を三重点 bundle と合わせ，
さらに cut の位置を離して選ぶことで，closed corner と切詰めた neck の
全ての実際の交わりで fiber partition を一致させる．
本節の結論は部分 fibration の接着である．Core への延長と全域被覆は
定理\ref{ke:actual-global-fibration}，全曲面を通る経路の距離比較は
定理\ref{ke:global-ambient-lower}，\ref{ke:global-ambient-upper}で別に扱う．

\subsection{残余基底と end 座標の区別}

二重線 $D=D_{ij}$ と，節\ref{ke:neck-corner-transition-section}で固定した
許容射影を用いる．三重点の集合を $Q\subset D$ とする．
射影 Morse 変換の後の方程式は，束の局所 frame において
\[
 x'y'=t h_t(p),\qquad p\in D\setminus Q
\]
である．Neck の残余基底は元の $p\in D$ である．
$h_t$ は束の section を frame で表示した関数であり，
全基底上の一対一座標ではない．単純 source 近傍で $h_t'(p)\ne0$ のときに限り，
$z=h_t(p)$ を局所座標として $x'y'=tz$ を使う．

一つの三重点 $q_k$ の近傍では local base coordinate を $p_k$ とし，
\eqref{ke:eq:arm-gauged-coordinates}の第三座標を
\[
 Z_{k,t}(p_k)=\Xi_{kk}(p_k)=C_D(p_k,t)^2/h_{k,t}(p_k)
\]
と書く．添字 $k$ は $h_t$ の当該局所 frame での表示も表す．
\[
 Z_{k,0}(p_k)=p_k/C_D(p_k,0)
\]
は $q_k$ で単純零点を持つ正則座標である．
各固定 punctured annulus 上では $Z_{k,t}\to Z_{k,0}$ が全固定階で成り立つ．
従って $t$ を十分小さくすれば，同じ annulus の少し小さい部分で
$Z_{k,t}$ も一対一の正則座標である．これは $h_t$ を global coordinate とする主張ではない．

Global neck の frame と各 end の frame の変換は，残余基底だけに依存する．
補題\ref{ke:projected-source-normal-form}の臨界断面と恒等近傍の平方根は
frame 変更に関して同変なので，normal circle は両表示で厳密に同じである．
一つの三重点を無限遠とする base coordinate $p_g$ と束の frame を取れば，
全ての compact exhaustion 上で $\Omega_N=d\log x'_g\wedge dp_g$ を使える．
Local frame における同じ体積形式との差は base-only unit である．
Canonical primitive の差は normal generator を消す閉形式なので，
normal moment の加法定数は変わらない．
第二 period については同じ結論を主張しない．

\begin{lemma}[固定残余 compact 上で先に選べる normal cutoff]
\label{ke:fixed-stage-source-neck}
$P\Subset\operatorname{int}P^+\subset D\setminus Q$ を固定し，
source は $P$ の内部に，$P^+\setminus\operatorname{int}P$ は source の外に置く．
固定 frame で $r=e^{-n/2}$ とし，
\[
 U_{t,n}=\{p\in\operatorname{int}P^+,\ |x'_g|,|y'_g|<r\},
 \qquad
 K_{t,n}=\{p\in P,\ |x'_g|,|y'_g|\le r^2\}
\]
とする．$P,P^+$ に依存する定数 $C_P,\delta_P,c_P>0$ と $N_P$ があり，
任意の固定した $n\ge N_P$ に対して，$|t|$ を十分小さくすれば次が成立する．
正則 circle generator $V$ は outer の作用不変領域 $U_{t,n}$ 全体で
\begin{equation}\label{ke:eq:source-fixed-stage-generator-bound}
 \sup_{U_{t,n}}|V|_{g_t}\le C_P/n
\end{equation}
を満たす．従って各 orbit の長さも $C_P/n$ 以下である．
さらに $K_{t,n}$ と外側の境界の actual buffer は
$\delta_P$ 以上であり，内側の二つの normal cut の canonical moment は逆符号で
絶対値が $c_P/n$ 以上である．
従って $n$ をさらに大きくして $C_P/n\ll\delta_P$ とすれば，
centered normal Moser map $\phi_{N,t}$ の source と像は共に $K_{t,n}$ を含む．
その actual 変位は $C_P/n$ 以下である．
十分小さい対称 interval を $(-c_P/n,c_P/n)$ 内に選ぶと，
円作用の商は $(\mu_N,p)$ によりその interval と $P$ の積に位相同定される．
\end{lemma}

\begin{proof}
定理\ref{ke:projected-source-neck}の証明を，$r$ を先に固定して適用する．
固定された source-free 部分には定理\ref{ke:fixed-base-neck-moser}の比較を使い，
source 近傍には系\ref{ke:source-short-circle}を使う．有限被覆により
生成場の\emph{点ごとの}上界の定数を一つの $C_P$ にできる．
実際，source-free patch では
定理\ref{ke:fixed-base-neck-moser}の比較から
$|V|_{g_t}\le C/\min(s,M_t-s)$ であり，
source patch では系\ref{ke:source-short-circle}から
$|V|_{g_t}\le C/\min(s_x,s_y,L_t)$ を得る．
$U_{t,n}$ 上ではいずれの分母も $n-O_P(1)$ 以上で，
固定 frame の有界なずれは $n$ を大きくして吸収できる．
従って\eqref{ke:eq:source-fixed-stage-generator-bound}を得る．
これは単に一周の積分長から点ごとの速度を推論したものではない．
残余方向では $P^+\setminus P$ の固定幅に射影下界を使う．
Normal 方向では $r^2\le|x'_g|\le r$ または対称な shell の成分極限と
cusp 下界を使う．この shell の深度比は $2$ なので，$n$ を大きくした後にも
固定された $P^+$ 上で正の下界 $\delta_P$ を選べる．
同じ成分の fixed-shell canonical moment 評価は normal cuts 上で
$c_P/n$ の符号付き下界を与える．
これらは各固定 $n$ の後に $t$ を小さくする順序の評価である．

系\ref{ke:short-circle-image-buffer}により，共通 source と像を持つ centered flow を得る．
Source の固定点では生成場と chain primitive が零であり，flow は固定点を保つ．
Moment の radial 単調性と source の二軸の記述は
定理\ref{ke:projected-source-neck}と同じであり，商の位相同定を与える．
この段階では $P$ を動かした定数の一様性を主張していない．
\end{proof}

\subsection{全 stage に共通な reduced marking}

\begin{lemma}[一つの marked pants 木系の共通選択]
\label{ke:fixed-marked-tree-exhaustion}
固定した退化族の各二重線 $D=D_{ab}$ に対し，一つの embedded marked
simple-saddle foliation $\mathcal F_D^0$ を $D\setminus Q$ 上に選べる．
これは固定 compact $K_D$ の外で $|Z_{k,0}|=\mathrm{constant}$ の円を葉とし，
各 source は互いに異なる full regular circle leaves に属する．
全 critical leaves の小さい飽和 pants neighborhoods と，これらの source leaves の
小さい飽和 annuli は $K_D$ の内部に入る．

本節の全ての十分大きい exhaustion stages と，それぞれの閾値以下の全ての
複素 $t\ne0$ について，この一つの葉層を基準として neck を構成できる．
各 end の全 regular circle で切った compact reduction の Reeb tree は，
葉の labels $k\notin\{a,b\}$，saddle labels，および $\mu=0$ の source marks を保って，
一つの固定した $T_{ab}=T_{ba}$ に同定される．後続の
命題\ref{ke:link-relative-collar-gluing}の相対補間と
補題\ref{ke:whole-fiber-core-trim}の切詰めを行っても，この木系の組合せ型は変わらない．
これは一つの許容された木系の存在であり，任意に指定した木系又は作用周期への
適合を主張しない．
\end{lemma}

\begin{proof}
\emph{固定した pants と full leaves．}
三重点の互いに素な小円板を除いた球面を取り，一つの labelled boundary を root とする．
三つ穴球面を木に沿って貼り，各 pants に saddle 一個の標準 Morse function を置く．
Regular annuli 上で単調に接続し，saddle の値を互いに異なるよう選ぶと，
内部の critical points が $d-4$ 個の index-one points だけである first integral
$h_D^0$ を得る．各 end へ radial に延長する．
有限個の source 点の近傍で，boundary-relative な小さい diffeomorphism により
葉層を動かし，全 source 値を互いに異なる regular values にできる．
実際，各 source の互いに素な小円板内に，既に選んだ有限個の値と全 critical values を
避ける regular point を一つずつ取り，それを source 点へ移す bump isotopy を使えばよい．
その後の first integral を改めて $h_D^0$，その葉層を $\mathcal F_D^0$ とする．

Source 点 $q_r(0)$ を含む full regular circle を $L_r$ と書く．
各 $L_r$ の小さい飽和 annulus $A_r$ と，各 full critical leaf の飽和 pants
$Q_v$ を互いに素に取れる．これら全体と，その間の有限個の regular annuli を
一つの固定 compact $K_D$ の内部へ入れる．$K_D$ の外では end の radial 葉を使う．
End collars を長くしても新しい vertex は加わらず，その外端で切った Reeb tree は，
leaf edges の長さを変えるだけで同じ labelled $T_{ab}$ になる．
$D_{ab}=D_{ba}$ は同じ実際の曲線なので，二方向に別の pants 分解を選ばない．

\emph{Moving sources の一価な標識．}
補題\ref{ke:projected-source-normal-form}の $h_t$ は，$t=0$ を含む円板上で
正則である．同補題の臨界断面は陰関数定理による一意な断面であり，行列の平方根は
恒等行列の近傍の級数で定められている．従って $t^{1/2}$ の枝を選ぶ操作はない．
$h_0$ の単純零点 $q_r(0)$ ごとに，$\overline U_r\subset\operatorname{int}A_r$ となる
互いに素な座標円板 $U_r$ を固定する．正則陰関数定理と $\partial U_r$ 上の非消滅により，
ある一つの円板 $|t|<\epsilon_D$ 上で，各 $U_r$ 内の唯一の零点
$q_r(t)$ が正則かつ一価に定まる．特に source labels の置換は起こらない．
この選択は $K_D$ 上で先に行うので，後の exhaustion stage に依存しない．

$z_r(q_r(0))=0$ とし，$\chi_r$ を $0$ の近傍で $1$，$U_r$ の境界近傍で
$0$ となる固定 bump function とする．$a_r(t)=z_r(q_r(t))$ と置き，
\[
 z_r(S_t(p))=z_r(p)+\chi_r(z_r(p))a_r(t)\quad(p\in U_r),
 \qquad S_t(p)=p\quad(p\notin\textstyle\bigcup_r U_r)
\]
と定める．$|t|$ を小さくすれば微分は恒等写像に十分近く，境界近くで恒等なので
これは diffeomorphism である．$a_r(t)$ を $u a_r(t)$，$0\le u\le1$ と置き換えると
同じ支持を持つ isotopy を得る．従って $S_t\in\operatorname{Diff}_0(D\setminus Q)$，
$S_0=\mathrm{id}$ であり，実二変数 $(\operatorname{Re}t,\operatorname{Im}t)$ に滑らかに依存する．
全 $Q_v$ と boundary/end collars は点ごとに固定する．
また $S_t(A_r)=A_r$ であり，$S_t(L_r)$ は $q_r(t)$ を含み，
他の source を含まない full regular circle である．

\emph{全 stage の end 補正．}
$b_j$ を十分大きいところから始め，$K_D\subset\operatorname{int}P_j$ とし，
本節で使う end の補正 annuli を全て $K_D$ の外に置く．
命題\ref{ke:single-arm-reduced-graph}と
補題\ref{ke:reduced-collar-graph}で使う normalized coordinates では，
指定された reduced 葉は $a=Q_{j,\mu,t}(c,\theta)$ という小さい circle graph である．
同補題の明示的な写像
\[
 D_{j,\mu,t}(c,\theta)
 =\bigl(c+\chi(c)(Q_{j,\mu,t}(c,\theta)-c),\theta\bigr)
\]
を使い，外へ恒等写像で延長する．第一成分の $c$ 微分は正であり，
$\chi$ を $u\chi$ に変える isotopy も同じ支持を持つ．
Graph の小ささ，共通 moment interval，全 trajectories の余裕は，
次の定理\ref{ke:neck-corner-diagonal-matching}で選ぶ既存の stage 条件そのものである．
ここで新しい moving-base 一様評価は必要ない．

従って neck の reduced 葉族を
\[
 \mathcal F_{j,\mu,t}
   =(D_{j,\mu,t}\circ S_t)_*\mathcal F_D^0
\]
と選べる．二つの補正の支持は互いに素である．全 full critical leaves は
同じ $Q_v$ にあり，source leaves は全 $\mu$ で $S_t(L_r)$ のままである．
End の cut は $D_{j,\mu,t}$ による radial circle の像で行うので，
内側の葉を切断せず，その Reeb tree は同じ $T_{ab}$ である．
特に $\mu=0$ で一つの source だけを含む閉じた円葉を丸ごと持ち上げられる．
これは点の近傍の marking だけを保つ主張ではない．

\emph{Core collars と critical values．}
後続の相対接着では $q_a=\delta_j=4\pi/N_j$ の葉を anchor に取り，
Liouville flow による葉族との比較を
定理\ref{ke:saddle-moving-relative-interpolation}の
boundary-relative な diffeomorphisms によって行う．
各 moment の全 partition が一つの diffeomorphism の像なので，
full critical leaf の連結性と一 saddle の性質，regular annuli の接続順序が保たれる．
支持は $\mu=0$ を避けるため，source の full pinched fibers とその marking は変わらない．
Source 点 $q_r(t)$ を非零 moment の reduced 曲面上でも marked points として指定する．
Anchor では全て相異なる regular leaves 上にあるので，小さい飽和 annuli 内で
regularity と分離を保つよう band を縮め，同定理の fixed-marked-point 版を用いる．
全 fibers ごとの trim と cone filling は，この同じ labelled end を corner annuli で
接続するため，$N_j$ の変更によって別の木を選ぶ操作にはならない．

初期の $h_D^0$ では source 値と critical values を分離して選んだ．
その後も葉の partition は上の diffeomorphisms によって同定されるので，
異なる source leaves が合体したり，source が saddle の full leaf に入ったりしない．
別々の connected leaves に与える数値が，first integral の再表示により一致しても，
それだけでは quotient の点は同一視されない．
同じ connected leaf に二つの saddle が生じる変形は用いていない．
ここで固定した補助 first integral の値は，第二 action の値を指定する条件ではない．

\emph{Parameter を一周するときの範囲．}
$S_t$ は $t=0$ を含む円板上の一価な公式であり，
end では各固定 annulus 上の $Z_{k,t}/Z_{k,0}$ が恒等に近い正則 unit なので，
その対数も $t=0$ で $0$ となる枝で取れる．
Collar graph の補正は circle 座標 $\theta$ に関して周期的で，
labelled end を入れ替えない．各 $t$ でこの基準から marking を移すため，
$\arg t$ を一周させても source labels と抽象的な pants 木系を同じものに選べる．
これは $X_t$ 全体の parameter monodromy の自明化ではない．
K2 で共通にする周期 $1$ の normal circle とその向きは保つが，
torus の全 $H_1$ 基底を punctured parameter 円板全体で固定するとは主張しない．
第二 cycle は各固定 $t$ の共通 torus 上で選ぶ従来の範囲に留める．
相対 core 補間も各 $t$ で同じ anchor marking から行えばよく，
stage 間又は $t$ 全体で滑らかな fibration 族を作る必要はない．
有限個の二重線についてこれらの選択を同時に行い，木系 $T=(T_{ab})$ を得る．
\end{proof}

この補題は，既知の Morse 葉層の isotopy 一般論を新たな分類定理として主張するものではない．
Morse mapping space の path components の分類
\cite[Main Theorem]{Maksymenko}だけからは，指定した full critical leaves 又は
embedded pants marking の不変性は従わない．
ここでは小円板と互いに素な end annuli に支持を持つ具体的な isotopies を用いた．

\subsection{三つの尺度を分ける diagonal 選択}

\begin{theorem}[Source-inclusive neck と全 end の開 collar の一致]
\label{ke:neck-corner-diagonal-matching}
固定二重線 $D$ 上に，source を相異なる regular leaves に置く
marked simple-saddle foliation を取り，三重点の punctured ends では
$|Z_{k,0}|=\mathrm{const}$ の円を葉とする．
この foliation は，source と saddle を含む固定 compact 部分の外では radial とする．
任意に遅い段階的な選択として，列
\[
 b_j\longrightarrow\infty,\qquad
 n_j\longrightarrow\infty,\qquad
 \iota_j>0,\qquad t_j\downarrow0
\]
を次の性質を満たすよう選べる．
$t_{j+1}<|t|\le t_j$ のとき，$I_j=[-\iota_j,\iota_j]$ 上に
source と simple saddles を含む actual $\omega_t$-Lagrangian neck があり，
各 end の開 collar で三重点側の指定 normal-first torus bundle と
fiber partition が厳密に一致する．Normal Moser map と canonical normal moment も一致する．
Source の一重 pinched torus，saddle の $S^1\times\figeight$，
および内部の marking は保持される．

この一致 collar は，triple の射影座標で
\[
 b=-\log|Z_{k,t}(p_k)|^2\asymp b_j,
 \qquad a,c\ge4n_j
\]
となる moment strip から構成する．$n_j/b_j$ は十分大きく選び，
$b\asymp b_j$ で $a$ または $c$ が $b_j$ 程度しかない領域まで
neck が覆うとは主張しない．
法線円作用の Moser 変位は neck 全体で $1/j$ 以下にできる．
また $m_t=b_j/64$ と取った三重点の cut は
\[
 m_t\to\infty,\qquad \log m_t=o(R_t)
\]
を満たすように選べるので，定理\ref{ke:slow-cut-closed-corner-gh}と両立する．
この結論は時間 parameter に関する微分可能性を要求しない．
\end{theorem}

\begin{proof}
\emph{残余領域を最初に固定する．}
初期葉層と動く source の標識には，補題\ref{ke:fixed-marked-tree-exhaustion}の
一つの共通選択を用いる．
十分大きい $b_j\uparrow\infty$ を取り，各 end の
$s_{k,0}=-\log|Z_{k,0}|^2$ を用いて
\[
 P_j=D\setminus\bigcup_k\{s_{k,0}>8b_j\},\qquad
 P_j^+=D\setminus\bigcup_k\{s_{k,0}>16b_j\}
\]
とする．大きい $j$ では，source と saddle は全て両領域の固定内部にあり，
残余 annuli $b_j/2<s_{k,0}<4b_j$ は互いに素である．
これらは各段階では固定 compact 領域であり，
補題\ref{ke:fixed-stage-source-neck}の定数を
$C_j,\delta_j,c_j,N_{P,j}$ と書く．
$N_{P,j}$ はこの固定残余領域に対する normal cut の下限であり，
後で core を作るために選ぶ component cut $N_j$ とは別の量である．

各 end の local frame，$C_D$，global neck の frame の間に現れる
対数絶対値の差を，上の少し広い固定 annuli 上で $M_j$ 以下に抑える．
$t$ に依存するものも，固定領域で $t=0$ の近傍へ制限して同じ上界に入れられる．
従って local normal 深度と global frame の深度の差は $M_j$ 以下である．
無限遠に置いた end の frame 変換にもこの bound を適用する．

\emph{Normal cutoff を次に選ぶ．}
$n_j$ は増加列とし，
\[
 n_j\ge N_{P,j},\qquad n_j>4M_j,\qquad
 n_j>100(1+2\sigma)b_j,\qquad
 \frac{C_j}{n_j}<\min\left(\frac1j,\frac{\delta_j}{100}\right)
\]
を満たすよう選ぶ．必要なら最初の有限個の項は除く．
$r_j=e^{-n_j/2}$ とし，global inner tube の深度は $2n_j$，
outer tube の深度は $n_j$ となる．

End では，命題\ref{ke:arm-common-moment-strip}と同じ canonical volume の計算を，
固定比 $a=Kb$ の代わりに $a=4n_j$ で行う．
Dominant-product 比較は $a/b$ に定数が依存しないので，
$b_j/2<b<4b_j$ において
\begin{equation}\label{ke:eq:exhaustion-normal-cut-moment}
 \mu,\nu\big|_{a=4n_j}
 =\frac{\pi}{n_j}
  +O\left(\frac1{n_jb_j}
         +\frac1{n_j}\left(\frac{n_j+b_j}{L_t}\right)^2\right)
\end{equation}
となり，$c=4n_j$ では主項の符号が逆になる．
座標接続の指数誤差も $C/(n_jb_j)$ に吸収される．
この計算では小さい triple primitive を使い，
系\ref{ke:transition-normal-primitive}によって normal moment を global neck の
primitive と厳密に同定する．

\emph{Moment interval を選ぶ．}
各 $j$ について
\[
 0<\iota_j<\min\left\{\frac{c_j}{16n_j},\frac{\pi}{64n_j}\right\}
\]
とする．構成は目標 interval $I_j$ より広い $|\mu|<2\iota_j$ の上で行う．
固定段階の global normal cuts と
\eqref{ke:eq:exhaustion-normal-cut-moment}の四つの境界値の間に
この広い interval が入るよう，後で $t$ の閾値を取る．
Normal 単調性により，三重点側の $\nu$-level と $\mu$-level は全て
$a,c>4n_j$ にある．Global neck 側では，source を含む全 $P_j$ 上で
同じ interval の円，または $\mu=0$ の source 固定点を持つ．

\emph{最後に $t_j$ を選ぶ．}
各段階で有限個の要求を同時に満たす閾値を選ぶ．具体的には，固定 $P_j^+$ と
$r_j$ に対する補題\ref{ke:fixed-stage-source-neck}の全評価，
source 零点の内部への保持，frame と Morse 座標の共通性，
$|t|\sup_{P_j^+}|h_t|<r_j^8$，および
\eqref{ke:eq:exhaustion-normal-cut-moment}の誤差を $\pi/(4n_j)$ 以下にする条件を課す．
さらに各固定 annulus 上で $Z_{k,t}/Z_{k,0}$ の正則対数が存在し，
\[
 (q_t,\theta_t)=
 \left(\log\frac{-\log|Z_{k,t}|^2}{b_j},\arg Z_{k,t}\right)
\]
とその $t=0$ 版の $C^3$ 差が $b_j^{-2}$ 以下になるようにする．
この微分は固定 annulus の座標での微分であり，各 $j$ を固定した後には
$t\to0$ により実現できる．

$t_j<t_{j-1}/2$ とし，さらに
\[
 R_t\ge j\bigl(1+\log(n_j+1)\bigr),\qquad
 L_t>100n_j
 \quad\text{for }0<|t|\le t_j
\]
を課す．固定段階の定数や $n_j$ に大きな増大があっても，
これらは全て $t_j$ を小さくすることで満たせる．
各評価は $0<|t|\le t_j$ 全体で成り立つので，
$t_{j+1}<|t|\le t_j$ で第 $j$ 段階を使える．
この選択から $\log m_t/R_t\le1/j$ が従う．

\emph{全 flow の共通開領域を確保する．}
内側の残余 annulus $b_j<b<2b_j$ を取り，外側を
$b_j/2<b<4b_j$ とする．上の座標比較により，この外側 annulus は
全て $P_j$ の内部にある．Normal 深度 $a,c>3n_j$ の部分は，
global frame に移しても深度 $>3n_j-M_j>2n_j$ なので $K_{t,n_j}$ 内にある．
また $n_j\gg b_j$ により，外側 annulus 全体で固定比の正則 plateau arm 内にある．

$\nu$-level の全 torus は $a,c>4n_j$ にあり，第二段の $\psi_s$ の
モデル長さと $C^1$ 誤差は $Cb_j^{-2}$ 以下である．
同じ領域で normal generator の長さは $C/n_j$，模型との誤差は $C/b_j$ なので，
centered normal Moser 場のモデル長さは $C/(n_jb_j)$ 以下である．
正規化 log 深度と残余 $q$ の微分は一様に有界であるから，
大きい $j$ では，第二段の全画像と，そこからの第一段の全 trajectories は
$a,c>3n_j$ と外側残余 annulus の中に留まる．
逆向きの flow と平均順序比較の小さい trajectories も，少し広い同じ margin に収められる．
比較的 compact な飽和開近傍を取れば，
全円 orbit，residual-stage の全画像，および normal-stage の全 trajectories を含む
共通の開 plateau が得られる．

Global neck と三重点側で同じ centered kernel を用いる．
両側の bare normal circle は正則座標同定 $G$ により厳密に共役されるので，
補題\ref{ke:normal-flow-naturality}から
\[
 \phi_{N,t,s}\circ G=G\circ\phi_{\alpha,t,s}
\]
がこの共通領域で成り立つ．Primitive の近似だけから flow の一致を推論していない．
共通 canonical moment は上で確認したものである．

\emph{Reduced 葉を合わせる．}
命題\ref{ke:single-arm-reduced-graph}をより狭い $I_j$ に制限する．
三重点 torus の第二段の像は，neck の同じ $\mu_N$-level の reduction 上で
$(q_t,\theta_t)$ に関して $C^1$ 誤差 $O(b_j^{-2})$ の円 graph を与える．
固定座標 $(q_0,\theta_0)$ への変更も上で $O(b_j^{-2})$ に抑えたので，
これらの円は基準の radial 葉の小さい graph である．
補題\ref{ke:reduced-collar-graph}を互いに素な広い annuli で用い，
各 moment に滑らかに依存し，内部の source と saddle 近傍では恒等な
foliation の変形を作る．無重みの moment 微分に一様な小ささは課さない．
命題\ref{ke:moment-dependent-neck-foliation}によって，
この葉層族を全 neck へ持ち上げられる．
Source の $\mu=0$ を含む全商では $(\mu_N,p)$ の位相同定を用い，
滑らかな reduction 座標を使うのは source-free な end のみである．

系\ref{ke:corner-leaves-in-neck-reduction}により，選んだ end の torus は集合として
\[
 G\phi_\alpha\psi(T)=\phi_N G\psi(T)
\]
であり，変形した neck 葉の持ち上げと完全に一致する．
三重点の三つの端に別々の normal-first 順序が必要な場合は，
定理\ref{ke:relative-corner-patch}を使う．支持は互いに素な arm にあり，
本証明で確保した共通 source/image の余裕の中に選べる．
これで一つの corner bundle の指定された各 end strip と一致させられる．

最後に，neck の終端を一致 annulus 内の一つの regular 葉で切り詰める．
例えば基準座標で $q_0=\log(3/2)$ の円を用い，その moment-dependent な像を
境界として取る．これは foliation に飽和した切断であり，
その両側に一致する開 collar がある．内部の sources と saddles は変形の支持外にあるので，
定理\ref{ke:projected-source-neck}の特異 fiber と marking はそのまま保持される．
この切詰めにより，離れた部分で二つの bundle の偶然の交差があることを
そのまま接着と呼ぶ必要もない．
\end{proof}

この議論は有限個の二重線について同時に行える．
残余 cut を共通にし，必要な $n_j$ の最大値，moment 幅と $t_j$ の最小値を取ればよい．
各三重点での三 arm の座標は既に一つの共通 chart に合わせている．
ただし，これで得るのは指定された end strips の相対的一致である．
浅い normal 領域と component core を含む全曲面の被覆は，
補題\ref{ke:fixed-stage-stratum-coverage}と定理\ref{ke:actual-global-fibration}で確認する．
Source を含む全 fiber 合併の径は命題\ref{ke:collapsed-tree-full-fiber-collapse}，
全曲面の no-shortcut は定理\ref{ke:global-ambient-lower}で評価する．
これらは局所 flow の小さい変位から直接従うものではない．

\subsection{Closed corner と切詰めた neck の全交わり}

Germ の一致から二つの実際の領域の和集合上の fibration を得るには，
指定した collar 以外の交点も除外する必要がある．そのため normal-first の相対化を，
最終的に使う moment strip より先に，少し浅い normal 領域まで施す．
以下の証明では，交点の moment を同定した後に初めて深い normal 下界を使う．

\begin{corollary}[全交わりが共通の葉で飽和した部分 fibration]
\label{ke:neck-corner-saturated-union}
定理\ref{ke:neck-corner-diagonal-matching}の有限個の二重線と三重点について，
diagonal 選択と cut の位置を取り直せば，次を同時に実現できる．
三重点側では
\[
 m_t=6b_j,\qquad t_{j+1}<|t|\le t_j
\]
とした全 softmax domain $\mathcal H_{t,m_t}$ を使い，
三つの arm の normal-first 構成を相対化した actual corner bundle を
$\mathcal C_{t,J}$ とする．Neck の各 end は，対応する $\kappa$-torus の残余深度
$b=7b_j$ によって指定される moment-dependent reduced 葉で切り詰める．
その actual 像を $\mathcal N_{t,D}$ とする．

各 $\mathcal C_{t,J}\cap\mathcal N_{t,D}$ は双方の全 torus fibers の和であり，
そこで二つの fiber partition と向き付き normal circle は厳密に一致する．
End の第二 cycle はこの共通 torus 同定で移送して標識する．独立に先指定した
全整数基底同士の一致や第二 period の一致を主張するものではない．
各 end には空でない開 overlap がある．小さい開近傍を用いて同じ構成を行えるので，
これらの和集合上に一つの actual $\omega_t$-Lagrangian 部分 fibration が得られる．
Neck 内部の sources と simple saddles の局所モデルは変わらない．

この選択も $m_t\to\infty$，$\log m_t=o(R_t)$ を満たす．
従って各 closed corner について，定理\ref{ke:slow-cut-closed-corner-gh}の
内在・protected GH 評価は保持される．和集合が曲面全体を覆うこと，
component core への延長，およびこの和集合の外を通る shortcut の評価は，
本系の結論には含まれない．
\end{corollary}

\begin{proof}
\emph{定数と cut の位置を選ぶ．}
前定理の残余領域を
\[
 P_j=D\setminus\bigcup_k\{s_{k,0}>12b_j\},\qquad
 P_j^+=D\setminus\bigcup_k\{s_{k,0}>16b_j\}
\]
に替える．Frame 差 $M_j$ は少なくとも $3b_j<s_{k,0}<12b_j$ を含む固定 annuli で取り，
$n_j>4M_j$，$n_j>1000(1+2\sigma)b_j$ とする．
固定段階の全 buffer 条件，$C_j/n_j<1/j$，および有限個の outer neck tubes の
分離も，必要なら $n_j$ をさらに大きくして満たす．
実際，異なる二重線上の固定 compact $P_j^+$ は三重点を除いているので互いに素であり，
まずその十分細い束近傍を互いに素に選び，最後に $t$ を小さくすればよい．
各三重点 $q_J$ の固定 chart の閉包は，他の三重点と，
$D_{ab}$，$\{a,b\}\not\subset J$ なる非 incident な二重線からも離して選ぶ．
一般位置条件により $q_J$ を通る二重線は $J$ 内の三本だけである．
固定段階の outer neck tubes を細くし，最後に $t_j$ を小さくすれば，
非 incident な $\mathcal C_{t,J}$ と $\mathcal N_{t,D_{ab}}$ は交わらない．
従って以下では対応する end の交点だけを扱えばよい．
Moment 幅は前定理と同じく
\[
 0<\iota_j<\min\{c_j/(16n_j),\pi/(64n_j)\}
\]
とし，$|\mu|<2\iota_j$ 上で構成する．
式\eqref{ke:eq:exhaustion-normal-cut-moment}は残余 annulus を
$4b_j<b<10b_j$ に替えても成立する．
$t_j$ を最後に小さくして，この annulus 上の全 $\mu,\nu$-levels が
$a,c>4n_j$ にあること，および $P_j^+$ の各 end 上で
$Z_{k,t}/Z_{k,0}$ の必要な $C^3$ 比較を確保する．
ここで $a,c$ は $\Xi_k$ の normal 深度，$b=-\log|Z_{k,t}|^2$ である．

\emph{Moment の制限より広い領域で先に平均順序を合わせる．}
新しい triple 座標の torus 飽和領域として
\[
 E_k^-=\{5b_j\le b\le9b_j,\quad a,c\ge n_j/2\}
\]
を取り，その少し広い $E_k^0$ と，さらに広い
\[
 E_k^+\Subset\{4b_j<b<10b_j,\quad a,c>n_j/4\}
\]
を取る．最初の不等式には必要に応じて小さい開近傍を付ける．
例えば残余の中間 cut を $4.5b_j,9.5b_j$，normal の中間 cut を $2n_j/5$
に取れば，入れ子領域の buffer は正の定数で下から抑えられる．
$\log(b/b_j),\log(a/n_j),\log(c/n_j)$ の滑らかな関数で作る
cutoffs の固定階の $h$-norms も一様である．
三つの $E_k^+$ は互いに素で，全て最小深度 $m_t/2=3b_j$ より深い．

座標接続の逆写像を $G$ とし，$\alpha=G^*\omega_t$ に対して，
一段階 map $\Phi$ も各 normal-first map $\phi_{\alpha,k}\psi_k$ も
centered kernel を用いて選ぶ．
定理\ref{ke:relative-corner-patch}を上の領域へ適用し，
\[
 C_t=G\Upsilon\Phi,\qquad
 C_t(z)=G\phi_{\alpha,k}\psi_k(z)\quad(z\in E_k^-)
\]
を得る．$C_t(\mathcal H_{t,m_t})=\mathcal C_{t,J}$ と置く．
$\Upsilon\Phi$ の変位と一次微分の変化は $O(b_j^{-2})$，
$G$ は指数的小摂動である．
従って，source と像のどちらでも各正規化 log 深度は $O(b_j^{-2})$ しか変わらない．
この相対化は moment を $I_j$ に制限する前に行っている．

\emph{Neck を共通の reduced 葉で切る．}
$4b_j<b<10b_j$ と $|\mu|<2\iota_j$ の共通 strip に
命題\ref{ke:single-arm-reduced-graph}を適用する．
$B$ を source の $\kappa$-torus の残余深度とし，
その第二段の像の reduced 葉を $\gamma_{\mu,B}$ と書く．
$B$ は torus 全体で一定であり，像の点ごとの残余深度 $b$ と区別する．
Graph 評価から
\[
 b=B\bigl(1+O(b_j^{-2})\bigr)
 \quad\text{on }\gamma_{\mu,B},
\]
またこれらの円は $B$ の順序を保つ．
補題\ref{ke:reduced-collar-graph}による foliation の移送を，
$5b_j<B<9b_j$ 上ではこの全 family と一致し，
source・saddle の近傍では恒等となるように取る．
Neck source を $\gamma_{\mu,7b_j}$ で切り，残余内部を残す．
その内側では，end 座標が定義される範囲で
$b\le7b_j+O(b_j^{-1})$ である．
実際の neck はこの領域を global normal Moser map $\phi_N$ で移したものとし，
その moment 座標は
\[
 \widehat\mu_N=\mu_N\circ\phi_N^{-1}
\]
である．平均前の点における $\mu_N$ と混同しない．

\emph{Neck の全 normal path の残余方向を閉じ込める．}
Global normal flow の全 trajectories は outer tube 内にあり，
global frame での二つの normal 深度は $n_j$ より大きい．
$4b_j<b<10b_j$ にいる間は frame 差から local normal 深度は
$n_j-M_j>3n_j/4$ である．
Projected 深度から旧座標の arm 条件へ戻す粗評価も確認する．
補題\ref{ke:arm-coordinate-ratio}の scaled Morse 写像の逆写像は，
$|x'|+|y'|<c|p|$，$|t|<c|p|^3$ 上で
\[
 |x|+|y|\le C\bigl(|x'|+|y'|+|t|/|p|^2\bigr)
\]
を満たす．これは臨界断面 $a_t=O(t/p^2)$，行列平方根
$S-I=O(t/p^3)$ と固定幅の逆写像定理から従う．
現在の固定 annulus を含む $A_j\Subset D\setminus Q$ 上で，
最後の $t_j$ の選択に
\[
 |t|\sup_{A_j}|p|^{-2}\le e^{-n_j/2},\qquad
 |t|\sup_{A_j}|p|^{-3}<c
\]
を追加する．固定 unit の上下界と $n_j\gg b_j$ により逆写像の領域条件が満たされ，
$|x'|,|y'|\le C e^{-3n_j/8}$ から旧 normal 深度は
$3n_j/4-O(1)$ 以上となる．$q_k=p+\alpha x+\beta y$ と
$Z_{k,t}/Z_{k,0}\to1$ から旧 residual 深度は $b+O(1)$ である．
従って旧最小深度は $3b_j$ より大きく，旧 normal/residual 深度比は
$1+2\sigma$ より大きい．この評価は projected normal circle 全体で保たれるので，
座標補間の cutoff は全 orbit 上で $\chi_k=1$ となる．
よってその全 normal circle は正則な $\Xi_k$ plateau に含まれる．
補題\ref{ke:normal-flow-naturality}をここでは vector field の段階で用いると，
$q=\log(b/b_j)$ に対して
\begin{equation}\label{ke:eq:neck-residual-first-exit}
 |Y_N|_h\le\frac{C}{n_jb_j},\qquad
 |dq(Y_N)|\le\frac{C}{n_jb_j}
\end{equation}
を得る．この評価には global と local の endpoint maps の一致は使わない．
Centered homotopy が不変模型を消すこと，normal generator の長さが $C/n_j$ 以下で
actual/model 誤差が $C/b_j$ 以下であることから直接従う．

Source の residual cut と $8b_j$ の間には固定の $q$-幅がある．
従って最初に $b=8b_j$ に到達する時刻があると仮定すると，
その直前の annulus を通る $q$-変位が
\eqref{ke:eq:neck-residual-first-exit}の時間 $1$ の積分より大きくなり矛盾する．
よって neck の全 image と全途中像は $b<8b_j$ にある．
Source が end の外にある場合も，同じ annulus を横切る必要があるので同じ議論になる．

\emph{実際の交点を先に広い plateau に入れる．}
$x\in\mathcal C_{t,J}\cap\mathcal N_{t,D}$ とする．
Corner の source は全深度 $\ell_a>6b_j$ を持つので，
小さい変位により $x$ の旧 triple 深度は
$\ell_a(x)\ge6b_j-O(b_j^{-1})$ である．
元の射影座標 $p$ は固定 affine triple 座標 $q_i,q_j,q_k$ の線形式であり，
それらと $\zeta_a$ の比は固定 unit である．従って
$|p(x)|\le C e^{-3b_j}$，すなわち
\[
 b(x)\ge6b_j-C>\tfrac{11}{2}b_j
\]
となる．ここでは $F=\Xi_k$ をまだ用いず，
固定 end 上の $Z_{k,t}/Z_{k,0}$ の比較だけを用いている．
既に得た neck の上界と合わせて
$11b_j/2<b(x)<8b_j$ を得る．

$x$ から global neck の逆 normal trajectory を時刻 $0$ まで追う．
全軌道は outer tube 内にあり，帯内では local normal 深度が $3n_j/4$ より大きい．
\eqref{ke:eq:neck-residual-first-exit}と固定幅の margin によって，
逆軌道も全時刻で $5b_j<b<9b_j$ に留まる．
従って全軌道と全 normal orbits 上で $F=\Xi_k$ が成立する．
特に $F(x)$ の残余深度は実際の $b(x)$ と同じである．

$x=C_t(z)$ と書く．$\Upsilon\Phi$ の小変位から，
$z$ の残余深度は $(5b_j,9b_j)$ 内，normal 深度は $n_j/2$ より大きい．
従って $z\in E_k^-$ であり，
$x=G\phi_{\alpha,k}\psi_k(z)$ という staged 表示を使える．
Global neck の逆軌道を $F$ で移すと local normal の逆方程式を満たすので，
ODE の一意性から
\begin{equation}\label{ke:eq:intersection-inverse-normal}
 \phi_N^{-1}(x)=G\psi_k(z)
\end{equation}
を得る．この時点までは，$z$ の normal 深度が $4n_j$ より大きいことを
仮定していない．

\emph{Moment と全 fiber を同定する．}
式\eqref{ke:eq:intersection-inverse-normal}，
系\ref{ke:transition-normal-primitive}，および第二段の moment の保存から
\[
 \widehat\mu_N(x)=\mu_N(G\psi_k(z))
                =\mu_\alpha(\psi_k(z))=\nu_k(z)
\]
が厳密に成立する．
$x$ が neck の目標 moment interval にあるので $\nu_k(z)\in I_j$ であり，
ここで初めて normal 単調性を使って $a(z),c(z)>4n_j$ と結論する．

$z$ を含む全 $\kappa$-torus を $T_z$ とし，その残余深度を $B$ とする．
Reduced graph の同定から，$p(G\psi_k(T_z))$ は
$\gamma_{\widehat\mu_N(x),B}$ そのものである．
Neck は $B=7b_j$ の葉で切られており，graph family は $B$ の順序を保つから，
交点が存在すれば $B\le7b_j$ である．
その全 torus 上で第二段と第一段の trajectories は
$a,c>3n_j$ と $4b_j<b<10b_j$ にあり，global inner tube 内に含まれる．
従って全 torus に自然性を適用でき，
\[
 C_t(T_z)=G\phi_{\alpha,k}\psi_k(T_z)
          =\phi_NG\psi_k(T_z)
\]
は neck の一つの全 fiber と一致する．
$\mathcal H_{t,m_t}$ は全 phase tori の和なので，
この torus の一部だけが corner に含まれることもない．
これで交わりが双方で飽和していることと fiber partition の一致が従う．

$|\mu|<\iota_j$ と，例えば $6.4b_j<B<6.6b_j$ の共通 levels は空でない．
Normal 深度は $4n_j$ より大きく，
$\min_a(\ell_a-m_t)\ge0.4b_j>1$ なので softmax の全比条件を厳密に満たす．
従ってこれは両側の内部に含まれる開 overlap である．
Closed corner の source を，全最小深度が $23b_j/4$ より大きい
小さい torus 飽和開近傍へ広げても，上の交点の下界と first-exit の余裕は保たれる．
Neck 側でも少し広い moment interval を使える．
こうして開領域の fiber partitions が接着する．
その overlap の基底座標は $(\nu_k,B)=(\widehat\mu_N,B)$ であり，
normal 単調性と reduced graph の順序保存により座標同定は滑らかである．
まず閉領域の基底を用いると，各基底は compact Hausdorff であり，
飽和した交わりの像は閉である．同じ全 fiber による同定は，proper quotient を介して
これらの閉集合の間の同相写像を与える．同定の graph は閉であり，
異なる corner 同士と異なる neck 同士は交わらないので，同値類の元は高々二つである．
従って有限 compact 空間上の同値関係は閉で，貼合せ基底は Hausdorff となる．
上の小さい開 collar を同じ $(\nu_k,B)$ 座標で延長すれば，
開領域の部分 fibration もこの基底の近傍上で得られる．

向き付き normal cycle は作用の同変性により一致する．第二 cycle は共通 torus
同定で移送する．Base-only frame unit の winding により，別に選んだ residual cycle
とは normal cycle の整数倍だけ異なる場合があるので，その二つの基底を同一視しない．
Source・saddle の内部標識は end の補正の支持外にあるため保たれる．

異なる三重点 chart は互いに素に選ぶ．異なる neck images は，最初に互いに素にした
outer tubes の中に留まる．従って有限個を同時に貼り合わせられる．
内部の singular fibers と marking は end の変形の支持外にある．

最後に $t_j$ の選択に
$R_t\ge j(1+\log(6b_j)+\log(n_j+1))$ を追加すれば，
$m_t=6b_j$ に対しても $\log m_t/R_t\le1/j$ である．
相対化の変位と $C^1$ 誤差は $O(b_j^{-2})$ なので，
節\ref{ke:relative-corner-section}の pullback metric の評価により
同じ closed-corner GH distortion bound が成立する．
この距離評価を neck，core，または保護領域の外まで延長してはいない．
\end{proof}
